\mysec{Lista 3}
\newcounter{tres}
\stepcounter{tres}
\myexe{\thetres}
\bigline
\textbf{Sejam $A, B \subseteq X$ conjuntos quaisquer e denote por $\chi_A$, $\chi_B$ suas respectivas 
funções características. Demonstre as seguintes propriedades:
\begin{enumerate}[label=(\alph*), wide, labelindent=0pt]
  \item $\chi_{A \cap B} = \chi_A\chi_B$
  \begin{proof}[\textbf{Dem:}]
    Note que, se $x \not \in A$ ou $x \not \in B$, segue que $x \not \in A\cap B$,
    desse modo, 
    $$\chi_{A\cap B}(x) = 0 = 0\cdot\chi_B(x) = (\chi_A\chi_B)(x)$$ ou 
    $$\chi_{A\cap B}(x) = 0 =\chi_A(x)\cdot 0 = (\chi_A\chi_B)(x).$$
    Caso contrário, se $x \in A$ e $x \in B$, 
    temos $x \in A\cap B$, portanto, $\chi_{A\cap B}(x)=1=1\cdot 1 = (\chi_A\chi_B)(x)$, como queríamos.
  \end{proof}
  \item $\chi_{A\cup B} = \chi_A + \chi_B - \chi_{A}\chi_B$
  \begin{proof}[\textbf{Dem:}]
  Note que, se $x \in A$ ou $x \in B$, temos $\chi_{A\cup B}(x) = 1$. Se $x$ está em ambos os conjuntos,
  $$\chi_A(x) + \chi_B(x) - (\chi_A\chi_B)(x) = 1 + 1 - 1 = 1 = \chi_{A\cup B}(x).$$ Se $x \not \in B$ 
  e $x \in A$, $\chi_A(x) +\chi_B(x) - (\chi_A\chi_B)(x) = 1 + 0 - 0 = 1 = \chi_{A\cup B}(x)$. De modo análogo segue 
  o caso em que $x \not \in A$ e $x \in B$. Para $x \not \in A\cup B$, temos $x \not \in A$ e $x \not \in B$, 
  ou seja, 
  $$\chi_{A\cup B}(x) = 0 = 0 + 0 - 0 = \chi_A(x) + \chi_B(x) - \chi_A\chi_B(x).$$
  Com isso, temos o que queríamos.
  \end{proof}
\end{enumerate}}
\biglinend 
\newpage
\stepcounter{tres}
\myexe{\thetres}
\bigline
\textbf{Sejam $(X, \mcal F, \mu)$ um espaço de medida e $(g_n)$ uma sequência em 
$\mcal M^+(X, \mcal F)$. Demonstre que: 
\begin{align*}
  \int\parth{\sum_{n=1}^\infty g_n}d\mu = \sum_{n=1}^\infty\parth{\int g_n d\mu}
\end{align*}}
\begin{proof}[\textbf{Dem:}] Das propriedades de integrais de funções
  em $\mcal M^+(X, \mcal F)$, segue que para todo $N \in \mbb N$, 
\begin{align*}
  \int\parth{\sum_{n=1}^N g_n}d\mu = \sum_{n=1}^N\parth{\int g_n d\mu}.
\end{align*}
  Visto que $f_N = \sum_{n=1}^N g_n$ é não-decrescente e 
  converge para $f = \sum_{n=1}^\infty g_n$, do Teorema da Convergência Monótona, 
\begin{align*}
  \int f d\mu = \lim_{N \rightarrow \infty} \int f_N d\mu \implies 
  \int \parth{\sum_{n=1}^\infty g_n} d\mu = \lim_{N \rightarrow \infty} 
  \int \parth{\sum_{n=1}^N g_n} d\mu.
\end{align*}
Desse modo, 
\begin{align*}
  \int \parth{\sum_{n=1}^\infty g_n} d\mu = \lim_{N \rightarrow \infty} 
  \int\parth{\sum_{n=1}^N g_n}d\mu = \lim_{N \rightarrow \infty }
  \sum_{n=1}^N\parth{\int g_n d\mu} = \sum_{n=1}^\infty\parth{\int g_n d\mu},
\end{align*}
como queríamos.
\end{proof}
\biglinend 
\newpage
\stepcounter{tres}
\myexe{\thetres}
\bigline 
\textbf{Sejam $(X, \mcal F, \mu)$ um espaço de medida e $\varphi \in \mcal M^+(X, \mcal F)$
uma função simples com representação (não necessariamente padrão)
\begin{align*}
  \varphi = \sum_{k =1}^mb_k \chi_{F_k},
\end{align*}
em que $b_k \in \bR$ e $F_k \in \mcal F$, para todo $k=1, \cdots, m$. Mostre que 
\begin{align*}
  \int \varphi d\mu = \sum_{k=1}^m b_k \mu(F_k).
\end{align*}}
\begin{proof}[\textbf{Dem:}]
  Note que, para cada $x \in X$, temos $\varphi(x) = \sum_{k \st x\in F_k} b_k$.
  Visto que $$\Im(\varphi)\subseteq \set{\sum_{k \in J} b_k \st J \subseteq \{1, \cdots, m\}},$$
  segue que existe uma sequência finita de valores distintos 
  $\{a_j\}_{j=1}^p = \Im(\varphi)$. Daqui, considere 
  uma sequência $\{E_j\}_{j=1}^p$ em $\mcal F$, tal que, 
  $E_j = \{x \in X \st \varphi(x) = a_j\}.$
  Desse modo, $\bigcup E_j = X$ e $\set{E_j}$ é disjunto pela definição de $\{a_j\}$.
  Sendo assim, como $\set{E_j}$ cobre $X$, cada $F_k$
  pode ser escrito como $F_k = \bigcup_{j = 1}^p (F_k \cap E_j)$. Visto que $E_j$ é disjunto,
  fixado $k \in \bN$, temos $\set{F_k \cap E_j}$ disjunto, ou seja, 
  $$\mu(F_k) = \mu\parth{\bigcup_{j=1}^p (F_k \cap E_j)} = \sum_{j=1}^p \mu(F_k \cap E_j).$$
  Portanto, 
  \begin{align*}
    \sum_{k=1}^m b_k \mu(F_k) = \sum_{k=1}^m b_k \sum_{j=1}^p \mu(F_k \cap E_j)
    =\sum_{j=1}^p \cparth{\sum_{k= 1}^m b_k\mu(F_k \cap E_j)}.
  \end{align*}
  Agora, note que pela definição de $\set{a_j}$, para cada $x \in E_j$,
  segue que 
    $\varphi(x) = a_j = \sum_{k=1}^m b_k \chi_{F_k}(x)$, 
    desse modo, 
    $$a_j\chi_{E_j} = \sum_{k=1}^m b_k \chi_{F_k}\chi_{E_j} 
    = \sum_{k=1}^m b_k \chi_{E_j \cap F_k}.$$
  Portanto, $a_j \mu(E_j) = \sum_{k=1}^m b_k\mu(E_j \cap F_k)$, e, já que 
  $\varphi = \sum_{j=1}^p a_j \chi_{E_j}$ é representação padrão de $\varphi$,
  \begin{align*}
    \sum_{k=1}^m b_k \mu(F_k) = \sum_{j=1}^p\cparth{\sum_{k=1}^m b_k \mu(F_k \cap E_j)}
    = \sum_{j=1}^p a_j\mu(E_j) = \int \varphi d\mu,
  \end{align*}
  como queríamos.
\end{proof}
\biglinend
\newpage
\stepcounter{tres}
\myexe{\thetres}
\bigline
\textbf{Considere o espaço de medida $(\bN, \mcal P(\bN), \mu)$ em que $\mu$
é a medida da contagem. Se $f$ é uma função não negativa em $\bN$, 
então $f \in \mcal M^+(\bN, \mcal P(\bN))$ e 
\begin{align*}
  \int f d\mu = \sum_{n=1}^\infty f(n).
\end{align*}}
\begin{proof}[\textbf{Dem:}]
  Seja $f: \bN \rightarrow \xbR$ não negativa. Visto que para todo $E \in \overline \borelR$,
  temos $\preim f (E) \in \mcal P (\bN)$, segue que $f \in \mcal M^+(\bN, \mcal P (\bN))$.
  Mais ainda, dado qualquer $\varphi \in \mcal M^+(\bN, \mcal P (\bN))$ simples 
  satisfazendo $0 \le \varphi \le f$ pontualmente, note que,
  $\varphi = \sum_{n \in \bN} \varphi(n)\chi_{\set{n}}$
  com $\varphi(n) \neq 0$ para finitos $n \in \bN$, é uma representação de $\varphi$. 
  Do exercício anterior,  $\int \varphi d\mu = \sum_{n \in \bN} \varphi(n)\mu(\set{n})$, 
  e, como $\mu$ é a medida da contagem, $\int \varphi d\mu = \sum_{n \in \bN} \varphi(n)$. 
  Daqui, como $\varphi \le f$ e ambos são não negativos,
  \begin{align*}
    \sum_{n \in \bN} \varphi = \int \varphi d\mu \le \sum_{n \in \bN} f,
  \end{align*}
  já que $\int f d\mu$ é o supremo de $\int \varphi d\mu$, segue que 
  \begin{align*}
    \int f d\mu \le \sum_{n \in \bN} f.
  \end{align*}
  Por outro lado, note que $f_N = \sum_{n = 1}^N f(n)\chi_{\set{n}}$ são funções simples 
  não negativas satisfazendo $0 \le f_N \le f$, para todo $N \in \bN$. Desse modo,
  \begin{align*}
    \sum_{n=1}^N f(n) = \int f_N d\mu \le \int f d\mu, \forall N \in \bN.
  \end{align*}
  Portanto, 
  \begin{align*}
    \lim_{N \rightarrow \infty} \sum_{n=1}^N f(n) = \sum_{n\in \bN} f(n) \le \int f d\mu.
  \end{align*}
  Com isso, temos o que queríamos.

\end{proof}
